\textbf{Phần 1. Trắc nghiệm lựa chọn}
\setcounterpageref{ex}{0}

\begin{ex} %[3D6N38]
Trong biểu thức chỉ có phép tính cộng, trừ thì ta thực hiện theo thứ tự
\choice
{Từ phải sang trái}
{Cộng trước, trừ sau}
{\True Từ trái sang phải}
{Tùy ý}
\loigiai{
Chỉ có cộng, trừ ta thực hiện từ trái sang phải.\\
\textbf{Đáp án của câu hỏi là C}
}
\end{ex}

\begin{ex} %[3D6N38]
Trong biểu thức có các phép tính cộng, trừ, nhân, chia thì thứ tự thực hiện là
\choice
{\True Nhân, chia trước rồi đến cộng, trừ}
{Cộng, trừ trước rồi đến nhân, chia}
{Từ phải sang trái}
{Không theo quy tắc}
\loigiai{
Nhân chia trước, cộng trừ sau.\\
\textbf{Đáp án của câu hỏi là A}
}
\end{ex}

\begin{ex} %[3D6H38]
Giá trị của biểu thức $15 + 25 - 10$ là
\choice
{$25$}
{$35$}
{$40$}
{\True $30$}
\loigiai{
$15 + 25 - 10 = 40 - 10 = 30$.\\
\textbf{Đáp án của câu hỏi là D}
}
\end{ex}

\begin{ex} %[3D6H38]
Giá trị của biểu thức $60 - 20 : 2$ là
\choice
{$20$}
{$40$}
{\True $50$}
{$30$}
\loigiai{
$60 - 20 : 2 = 60 - 10 = 50$.\\
\textbf{Đáp án của câu hỏi là C}
}
\end{ex}

\begin{ex} %[3D6H38]
Nếu trong biểu thức có dấu ngoặc $( )$, ta thực hiện
\choice
{Ngoài ngoặc trước}
{\True Trong ngoặc trước, ngoài ngoặc sau}
{Chỉ tính ngoài ngoặc}
{Bỏ ngoặc tùy ý}
\loigiai{
Thực hiện các phép tính trong ngoặc trước.\\
\textbf{Đáp án của câu hỏi là B}
}
\end{ex}

\begin{ex} %[3D6H38]
Giá trị của biểu thức $(45 - 15) \times 3$ là
\choice
{\True $90$}
{$80$}
{$100$}
{$75$}
\loigiai{
$(45 - 15) \times 3 = 30 \times 3 = 90$.\\
\textbf{Đáp án của câu hỏi là A}
}
\end{ex}

\begin{ex} %[3D6V38]
Tính giá trị biểu thức: $125 + 75 : 5$.
\choice
{$135$}
{\True $140$}
{$145$}
{$40$}
\loigiai{
$125 + 75 : 5 = 125 + 15 = 140$.\\
\textbf{Đáp án của câu hỏi là B}
}
\end{ex}

\begin{ex} %[3D6V38]
Tính giá trị biểu thức: $80 : (2 \times 4)$.
\choice
{$160$}
{$20$}
{$40$}
{\True $10$}
\loigiai{
$80 : (2 \times 4) = 80 : 8 = 10$.\\
\textbf{Đáp án của câu hỏi là D}
}
\end{ex}

\begin{ex} %[3D6V38]
Tính giá trị biểu thức: $240 - 40 \times 3$.
\choice
{$600$}
{\True $120$}
{$140$}
{$100$}
\loigiai{
$240 - 40 \times 3 = 240 - 120 = 120$.\\
\textbf{Đáp án của câu hỏi là B}
}
\end{ex}

\begin{ex} %[3D6V38]
Tính giá trị biểu thức: $(120 + 80) : 4$.
\choice
{$40$}
{$60$}
{$70$}
{\True $50$}
\loigiai{
$(120 + 80) : 4 = 200 : 4 = 50$.\\
\textbf{Đáp án của câu hỏi là D}
}
\end{ex}

\begin{ex} %[3D6V38]
Điền dấu thích hợp vào ô trống: $30 + 20 \times 2 \dots (30 + 20) \times 2$.
\choice
{\True $<$}
{$>$}
{$=$}
{Không so sánh được}
\loigiai{
Vế trái: $30 + 40 = 70$. Vế phải: $50 \times 2 = 100$. Do $70 < 100$ nên điền dấu $<$.\\
\textbf{Đáp án của câu hỏi là A}
}
\end{ex}

\begin{ex} %[3D6C38]
Tìm giá trị của $x$ biết: $(x + 25) \times 3 = 150$.
\choice
{$20$}
{$30$}
{\True $25$}
{$35$}
\loigiai{
$x + 25 = 150 : 3 = 50 \Rightarrow x = 50 - 25 = 25$.\\
\textbf{Đáp án của câu hỏi là C}
}
\end{ex}

\textbf{Phần 2. Trắc nghiệm đúng sai}
\setcounterpageref{ex}{0}

\begin{ex} %[3D6H38]
Xét tính đúng sai về giá trị của các biểu thức:
\choiceTF
{\True $100 - 50 + 20 = 70$.}
{\True $40 + 20 \times 2 = 80$.}
{$(40 + 20) \times 2 = 80$.}
{\True $(40 + 20) \times 2 = 120$.}
\loigiai{
\begin{itemize}
    \item $100 - 50 + 20 = 50 + 20 = 70$: Đúng.
    \item $40 + 40 = 80$: Đúng.
    \item $(40 + 20) \times 2 = 60 \times 2 = 120 \neq 80$. Ý 3 Sai, Ý 4 Đúng.
\end{itemize}
\begin{itemize}
    \item a) \textbf{Đúng}.
    \item b) \textbf{Đúng}.
    \item c) \textbf{Sai}.
    \item d) \textbf{Đúng}.
\end{itemize}
}
\end{ex}

\begin{ex} %[3D6V38]
Cho biểu thức $A = 120 : (3 + 3)$ và $B = 120 : 3 + 3$. Xét tính đúng sai:
\choiceTF
{\True Giá trị biểu thức $A$ là $20$.}
{\True Giá trị biểu thức $B$ là $43$.}
{Biểu thức $A$ có giá trị lớn hơn biểu thức $B$.}
{\True Biểu thức $A$ có giá trị bé hơn biểu thức $B$ ($20 < 43$).}
\loigiai{
\begin{itemize}
    \item $A = 120 : 6 = 20$. Đúng.
    \item $B = 40 + 3 = 43$. Đúng.
    \item $20 < 43$ nên $A < B$. Ý 3 Sai, Ý 4 Đúng.
\end{itemize}
\begin{itemize}
    \item a) \textbf{Đúng}.
    \item b) \textbf{Đúng}.
    \item c) \textbf{Sai}.
    \item d) \textbf{Đúng}.
\end{itemize}
}
\end{ex}

\textbf{Phần 3. Trả lời ngắn}
\setcounterpageref{ex}{0}

\begin{ex} %[3D6H38]
Tính giá trị biểu thức: $45 + 15 \times 3$.
\shortans{90}
\loigiai{
\textbf{Đáp số:} $90$\\
$45 + 45 = 90$.
}
\end{ex}

\begin{ex} %[3D6H38]
Tính giá trị biểu thức: $(60 - 20) : 4$.
\shortans{10}
\loigiai{
\textbf{Đáp số:} $10$\\
$40 : 4 = 10$.
}
\end{ex}

\begin{ex} %[3D6V38]
Tính giá trị biểu thức: $150 - 50 : 2$.
\shortans{125}
\loigiai{
\textbf{Đáp số:} $125$\\
$150 - 25 = 125$.
}
\end{ex}

\begin{ex} %[3D6V38]
Tính giá trị biểu thức: $(140 + 60) \times 3$.
\shortans{600}
\loigiai{
\textbf{Đáp số:} $600$\\
$200 \times 3 = 600$.
}
\end{ex}

\begin{ex} %[3D6V38]
Tính giá trị biểu thức: $300 : (2 \times 3)$.
\shortans{50}
\loigiai{
\textbf{Đáp số:} $50$\\
$300 : 6 = 50$.
}
\end{ex}

\begin{ex} %[3D6C38]
Tìm $x$ biết: $(x - 10) \times 4 = 160$.
\shortans{50}
\loigiai{
\textbf{Đáp số:} $50$\\
$x - 10 = 160 : 4 = 40 \Rightarrow x = 40 + 10 = 50$.
}
\end{ex}
